\section{目标函数与评价指标}
\label{sec:objective}

\subsection{评价范围}

评价指标以各模块的实际功率、状态和产品结算量为基础。运行调度考虑变动收入、变动成本、
弃电和服务违约；资本性支出、固定运维、融资、设备更换和税费在生命周期评价中另行核算。
因此，本章的运行净收益用于比较调度方案，不直接等同于项目利润、NPV或IRR。

\subsection{运行经济目标}

定义运行期经济净成本
\begin{equation}
\begin{aligned}
F_{eco}={}&C_{OM}+C_{start}+C_{deg}+C_{transport}
+C_{ENS}+C_{SLA}+C_{curt}+C_{terminal}\\
&-R_{elec}-R_{H2}-R_{compute}-R_{marine}.
\end{aligned}
\label{eq:economic-objective}
\end{equation}
各项统一以CNY计。电力收入按海缆受端电量电量计算：
\begin{equation}
R_{elec}=\sum_t\pi_t^{elec}P_t^{cab,r}\Delta t_t;
\label{eq:revenue-electricity}
\end{equation}
氢气收入按扣除通道损失后的接收量计算：
\begin{equation}
R_{H2}=\sum_t\pi_t^{H2}m_t^{H2,del};
\label{eq:revenue-hydrogen}
\end{equation}
算力和海洋服务收入分别依据4.6的实际服务量和式~\eqref{eq:marine-service}的实际供给量
计算。内部电力转移价仅用于分析业务单元边际贡献，在全系统目标中不重复计入。

$C_{deg}$按电池实际吞吐、SOC和温度相关寿命模型计量；$C_{start}$按实际增机或启动事件
计量；$C_{transport}$按管输/船运实际提取量或接收量计量；$C_{terminal}$用于反映库存
机会成本或滚动终端价值用于反映跨期价值，物理安全库存仍由状态约束保证。所有价格、违约、退化、运输和
终端价值参数须记录币值年、含税口径、适用区域和来源类型，在合同或企业成本数据明确前，相关参数采用基准值并在敏感性分析中给出取值范围。

绿氨尚未进入现行物理变量、资产账簿和输出与结算边界，故式~\eqref{eq:economic-objective}不
确认绿氨收入。待4.3补齐合成氨工艺后，应新增可追溯的产量、能耗、库存、运输和合同项。

\subsection{全生命周期经济性账簿}

对给定候选容量，资本回收因子为
\begin{equation}
CRF(r,n)=\frac{r(1+r)^n}{(1+r)^n-1},
\label{eq:crf}
\end{equation}
设备$i$的年化固定成本可写为
\begin{equation}
C_i^{annual}=CAPEX_i\,CRF(r,n_i)+C_i^{FOM}+C_i^{repl,annual}.
\label{eq:annualized-cost}
\end{equation}
该结构与公开技术经济评价中显式计入资本成本、寿命和固定运维的原则一致
\cite{nrel-atb,iec60300}。完整NPV按投运年现金流写为
\begin{equation}
NPV=-I_0+\sum_{y=1}^{Y}
\frac{R_y-C_y-I_y^{repl}}{(1+r)^y}.
\label{eq:npv}
\end{equation}
折现率、建设期支出曲线、税费、融资结构、寿命和更换年份需由第6章及企业财务边界
确认。当前财务参数用于方法展示，不构成项目投资评价结论。短时场景仅用于运行机制比较，不直接外推为正式NPV、LCOE、
LCOS或LCOH结论。

\subsection{环境目标}

环境目标采用固定边界下的项目排放减基准情景避免排放：
\begin{equation}
F_{env}=GHG_{project}-GHG_{avoided,baseline}.
\label{eq:environment-objective}
\end{equation}
运行代理可按源侧使用、电池吞吐、制氢、运输、氢转电和外购电等活动量乘相应排放因子，
并按预先设定的电力、氢气、算力和海洋服务基准计算避免排放。完整生命周期评价需依据
ISO~14040和ISO~14044完成目标与范围、清单、影响评价和解释
\cite{iso14040,iso14044}。设备BOM/EPD、海工船舶、备件、冷媒、替换、退役回收及多产品
分摊规则在取得企业或第三方LCA数据前采用统一的情景参数；现行代理仅用于场景比较，不作为已认证的
产品碳足迹。

\subsection{可靠性指标}

电力系统未供能仅统计具有物理功率需求的负荷：
\begin{equation}
ENS=\sum_t\left(P_t^{ENS,int}+P_t^{ENS,mar}
+P_t^{ENS,DC}\right)\Delta t_t,
\label{eq:ens}
\end{equation}
其中$P_t^{ENS,DC}$表示数据中心基础设施电力需求的未供功率。算力任务未完成量采用
MWh-CS计量，单独定义为
\begin{equation}
W^{CS,unserved}=\sum_tU_t^r+U^{\mathrm{term}},
\qquad
W^{CS,late}=\sum_tH_t^{\mathrm{late}}.
\label{eq:compute-reliability}
\end{equation}
电力ENS与算力服务缺口单位不同，分别报告，不进行直接相加。

当存在概率场景$s$及其概率$\pi_s$时，电力系统期望未供能为
\begin{equation}
EENS=\sum_{s\in\mathcal S^{scen}}\pi_sENS_s,
\qquad \sum_s\pi_s=1.
\label{eq:eens}
\end{equation}
EENS的定义和场景口径参考ENTSO-E资源充裕度方法\cite{entsoe-maf}。LOLE、LOLP、N-1、
频率最低点、ROCOF和黑启动成功率需要额外的概率或动态模型，不由本章小时模型直接给出。

\subsection{多目标集成}

由于CNY、kgCO$_2$e和MWh无法直接相加，主模型采用$\varepsilon$约束法：
\begin{equation}
\min F_{eco}
\quad\mathrm{s.t.}\quad
ENS\le\varepsilon_{ENS},\quad
W^{CS,unserved}\le\varepsilon_{CS},\quad
F_{env}\le\varepsilon_{GHG},\quad
\eta_{RE}\ge\varepsilon_{RE},
\label{eq:epsilon-constraint}
\end{equation}
其中可再生能源利用率为
\begin{equation}
\eta_{RE}=\frac{\sum_tP_t^{src,use}\Delta t_t}
{\sum_tP_t^{src,av}\Delta t_t}.
\label{eq:renewable-utilization}
\end{equation}
$\varepsilon$约束法参考Mavrotas的增强实现\cite{mavrotas2009}。可靠性和环境阈值根据
合同要求及方案比较目的设定；改变阈值可获得不同的经济性--可靠性折中方案。
